\chapter{Jacques Ellul and the Supremacy of Technique}
\label{ch:ellul}

The philosophical foundation of this monograph is taken from Jacques Ellul's \emph{The Technological Society}~\cite{ellul1964}. Ellul's central concept, technique, does not refer to machines. It refers to the totality of methods rationally arrived at and having absolute efficiency, in every field of human activity, for a given stage of development. A machine is one embodiment of technique. A bureaucratic procedure, an advertising campaign, and an accounting convention are others. What they share is that each is chosen because it is the most efficient available means.

The argument examined here is not that efficiency is undesirable. Ellul's claim is narrower and more demanding: that when efficiency becomes the supreme criterion, competing values are subordinated to it without ever being explicitly defeated in argument. A choice that is technically superior needs no further defense, and a choice that is technically inferior must justify itself against a standard it did not set. Over time the burden of proof shifts so completely that alternatives are rejected for failing a test that no one chose to apply.

\section{Efficiency and Legitimacy}

The chapter turns on a distinction between two propositions that ordinary speech runs together. The first is that an arrangement is efficient, meaning that it achieves a stated result at low cost. The second is that the arrangement is legitimate, meaning that it can be justified to the people it affects, with respect to ends they have reason to accept. The first is a technical claim, answerable to measurement. The second is a normative claim, answerable to argument. A very efficient arrangement can be illegitimate, as when it achieves a result at low cost by imposing the cost on parties who never consented. A legitimate arrangement can be inefficient, as when a community chooses a slower process because it values participation.

Technique, in Ellul's account, collapses the second proposition into the first. Because efficiency can be measured and legitimacy cannot, the measurable quantity comes to stand in for the whole. This is the point at which the argument of this monograph meets the argument of Chapter~\ref{ch:goodhart}. The elevation of efficiency to a supreme criterion is a special case of a measurable proxy acquiring authority over the purpose it was introduced to serve.

\section{Normative Authority Without Defense}

One of the more subtle observations in this tradition is that technical optimization can acquire normative authority without anyone defending the social arrangement it produces. No one need argue that a city should be organized around the automobile. The organization follows from a sequence of locally efficient decisions, each defensible, none of them a decision about the city. The aggregate arrangement is then treated as a given to which individuals must adapt.

Computational infrastructure exhibits the same pattern. A series of reasonable decisions to consolidate hardware, to raise utilization, and to amortize fixed costs over larger workloads yields an arrangement in which computation is concentrated in a small number of facilities and everyone else is a client. No single decision was a decision to centralize intelligence. The centralization is an emergent consequence, and its justification, where one is offered, is that it is the efficient outcome.

\section{Which Adapts to Which}

The question that organizes the rest of the monograph can now be stated. Should computational infrastructure determine the organization of human activity, or should it adapt to purposes established independently of it? The first answer treats the infrastructure as fixed and fits households, buildings, and institutions to it. The second treats the purposes as fixed, among them keeping a building warm, keeping food cold, and preserving a person's control over their own machine, and asks what computation can be done within them.

The second answer does not reject efficiency. It places efficiency in a subordinate position: among the ways of serving a purpose, choose the efficient one, but do not let the pursuit of efficiency redefine the purpose. This ordering is the sense in which the proposed architecture is a response to Ellul, and Chapter~\ref{ch:idle} returns to it as a principle of governance.

\section{Influence Without Instruction}

A closely related mechanism appears in an unrelated field and clarifies what is at stake. In a 2026 essay on technology journalism, Caffier~\cite{caffier2026} describes how corporate hospitality, including sponsored travel, managed tours, and gifts, can shape reporting without any request for favorable coverage. His point is that the reporter may truthfully state that no promises were made and that the sponsorship was disclosed, and still find that the access altered which questions seemed natural to ask and which stories seemed worth pursuing. A bribe tries to change a decision. Managed access changes the set of decisions that are psychologically and professionally available, and the stories that are never pursued do not appear as editorial choices that anyone could audit.

The mechanism can be stated compactly. Let $H$ denote the investigator's evidence, $A$ the access conditions established by the subject, and $Q$ the question the investigator pursues. Editorial independence guarantees that the subject cannot dictate the published conclusion. It does not guarantee
\[
P(Q\mid H,A)=P(Q\mid H),
\]
and the purpose of a sophisticated public relations operation may be precisely to make these distributions differ without issuing any instruction. Caffier's essay is a single first-person account, and its broader claims about the industry and about the mental states of the people he meets are not demonstrated by the encounters he describes. It is used here only to illustrate a mechanism, which parallels the one in this chapter: an optimized arrangement can acquire authority through the conditions it creates and not through any argument it makes.

Ellul's distinction is the useful one. An arrangement can be genuinely efficient at what it does, whether the arrangement is centralized hardware utilization or hospitality as a communication technique, and that achievement does not establish that the resulting distribution of knowledge or control is desirable. Even if the party optimizing were benevolent and its calculations entirely accurate, the underlying problem would remain. The question is not only whether the optimization works, but who decided that the thing being optimized should govern everything else.

\section{A Reactionary Objection}

The Colombian aphorist Nicolás Gómez Dávila supplies an objection that is more severe than Ellul's, and the monograph should answer it even though it does not require a chapter of its own. His scholia~\cite{gomezdavila1977}, brief observations written in a reactionary and Catholic mode, doubt that historical novelty, democratization, and growing technical capability constitute improvement, and the doubt extends to the modern faith in deliberate social reconstruction. His position should not be conflated with Ellul's. Ellul examines how efficiency, once supreme, displaces other forms of judgment. Gómez Dávila rejects the more general identification of the new with the better. Nor can he be recruited as a patron of distributed computing, since his skepticism of engineered social arrangements makes him more useful here as a critic of the project than as an authority for it. The objection below is constructed for the monograph from the stance he is known for, and it is not a claim about what he said of computers.

The objection is that an alternative technological system may reproduce the assumptions of the system it opposes in a different organizational form. The proposed architecture is itself designed, specified, and quantified. Its vocabulary of residual capacity treats the domestic sphere as a field of resources waiting to be organized, and it presupposes that computation discovered anywhere is worth seeking. A network of reused machines may look backward by conventional measures of performance and still be advantageous in resilience, autonomy, reversibility, and avoided capital. Those criteria are also measures, however, and they are exposed to the same corruption as the proxies of Chapter~\ref{ch:goodhart}. The distinction between progress as increasing capability and progress as preserving the conditions under which people can live well does not escape the difficulty by being drawn, because the second kind of progress can be turned into a metric and optimized just like the first.

The monograph's reply is partial. It does not claim to escape technique. It claims, in the remark that follows, that the order of priority between a purpose and its measure can be made explicit in the design of an institution. The architecture also carries limits on itself: the right to remain idle, opt-in participation, consent that expires, and stated conditions under which it fails. The one feature that answers the objection on its own terms is reversibility, since a system that can be abandoned without harm has not captured those who joined it. But reversibility asserted as an intention is not reversibility. Organizational commitments accumulate, as Chapter~\ref{ch:horizon} observes, and dependence grows through subsidies and bundled services, as Chapter~\ref{ch:idle} observes. A test is needed that does not rely on good intentions.

\begin{definition}[Dissolution condition]
\label{def:dissolution}
A network satisfies the \emph{dissolution condition} at time $t_0$ if its dissolution at $t_0$ would meet three requirements. First, functional restoration: each participating system performs its primary function no worse than before it joined. Second, financial restoration: every cost that a participant incurred solely because of participation, and has not recovered through the network's benefits, is borne by the network or its operator from a reserve held for the purpose. Third, release: no obligation to the network survives.
\end{definition}

The condition is demanding, and it is not met by default. A household that removed its conventional heater to install a computing module, or whose heating subsidy depends on continued participation, or whose hardware cannot be recovered, fails it. Read literally, full restoration is impossible, since participants spend time and attention that no reserve can return, and a participant who has come to rely on the network's services cannot be made as though the network had never existed. The condition is therefore stated in three parts that can be checked, and it asks for restoration of function and of identifiable costs, not of everything. It also separates two kinds of dependence. Dependence of a host's primary function on the network, as when a household loses heat when the network is disconnected, is prohibited, and the resistance fallback element of Chapter~\ref{ch:thermal} is the means of preventing it. Ordinary voluntary reliance on a service that exists only because the network does is permitted, since a network that provides nothing worth relying on has no case for existing.

The second requirement can be made quantitative. If $U_i(t)$ denotes the unrecovered participation cost of participant $i$ at time $t$, the network must hold a dissolution reserve $R(t)$ satisfying
\[
R(t)\ \ge\ \sum_i U_i(t),
\]
which places the cost of withdrawal on the network and not on its participants. The reserve is itself an organizational commitment of the kind identified in Chapter~\ref{ch:horizon}, and it reduces the network's attractiveness at small scale by the amount it ties up. Modular, retrievable hardware reduces each $U_i$ directly, and the no-conditionality rule of Chapter~\ref{ch:idle} keeps subsidy dependence from arising. A tension remains that no rule removes. The stronger the services of the network become, the more reasonably its participants will rely on them, and the greater the pressure to preserve the network for their sake. The condition cannot guarantee that nobody comes to depend on the network while the network also grows into substantial infrastructure. What it can do is keep dependence for essentials absent and make the cost of leaving explicit, so that the momentum of a useful network is not purchased with the captivity of its participants. The condition is necessary for the architecture to count as reversible and it is not sufficient for it to count as legitimate. A system can satisfy it and still reproduce the assumption that idle things exist to be organized, and the remainder of the objection is left open.

\begin{remark}
Ellul's pessimism about the autonomy of technique should be kept distinct from the engineering claims of later chapters. The proposal here is not that technique can be abolished or that its logic can be ignored. It is the more modest claim that the order of priority between a purpose and its measure can be made explicit in the design of an institution, and that doing so changes what the institution rewards.
\end{remark}
